\section{Diseño del filtro pasabanda para PCB}

La consigna del diseño propone un producto asociado al procesamiento de audio, cuyo objetivo es realizar un primer filtrado de la voz humana. Se pide una respuesta pasabanda de dos etapas, utilizando un TL082 o LM833 únicamente como buffer o como amplificador no inversor. Además, cada etapa puede ser de orden máximo dos.

\subsection{Especificación adoptada}

Para realizar un primer filtrado de voz se adoptó como banda útil aproximada:

\begin{equation}
    \boxed{\SI{300}{Hz}\leq f \leq \SI{3.4}{kHz}}.
\end{equation}

Esta banda corresponde a una zona de alta importancia para la inteligibilidad de la voz. Por lo tanto, el objetivo fue atenuar frecuencias menores a aproximadamente \SI{300}{Hz} y frecuencias mayores a aproximadamente \SI{3.4}{kHz}. La estrategia general fue combinar una etapa pasaaltos y una etapa pasabajos:

\begin{equation}
    H(s)=\Hpa(s)\Hpb(s).
\end{equation}

\subsection{Alternativa A: etapa RLC y etapa RC-RC}

La primera alternativa estudiada combinaba un pasabajos RLC con una etapa pasaaltos RC-RC, separados por un buffer. El objetivo era obtener el corte superior mediante el RLC y el corte inferior mediante el bloque RC.

\begin{figure}[H]
    \centering
\begin{tikzpicture}[transform shape]
	% Paths, nodes and wires:
	\draw (4.25, 1) to[sinusoidal current source, l={$v_i$}] (4.25, -1.25);
	\draw (4.25, 1) to[american resistor, l={$R_1$}] (6.25, 1);
	\draw (6.25, 1) to[american inductor, l={$L1$}] (8, 1);
	\draw (8, 1) to[capacitor, l={$C_1$}] (8, -1.25);
	\node[ground] at (8, -1.25){};
	\node[op amp, yscale=-1] at (10.333, 0.51){};
	\draw (8, 1) -| (9.143, 1.029);
	\draw (11.523, 0.51) to[capacitor, l={$C_3$}] (13, 0.5);
	\draw (13, 0.5) to[capacitor, l={$C_4$}] (15, 0.5);
	\draw (13, 0.5) to[american resistor, l={$R_3$}] (13, -1.25);
	\draw (15, 0.5) to[american resistor, l={$R_4$}] (15, -1.25);
	\draw (15, 0.5) -- (15.75, 0.5);
	\draw (9.143, 0.02) |- (11.5, -1.75) -| (11.5, 0.5);
	\node[ground] at (13, -1.25){};
	\node[ground] at (15, -1.25){};
	\node[ground] at (4.25, -1.25){};
    \node at (15.75, 0.5) [right] {$v_o$};
\end{tikzpicture}
    \caption{Alternativa A: pasabajos RLC, buffer y pasaaltos RC-RC.}
    \label{fig:opcion_a}
\end{figure}

Para el pasabajos RLC con salida en el capacitor:

\begin{equation}
    \Hpb(s)=\frac{1}{s^2L_1C_1+sR_1C_1+1}.
\end{equation}

Si se impone un corte superior cercano a \SI{3.4}{kHz} y se usa $L_1=\SI{1}{mH}$:

\begin{equation}
    C_1=\frac{1}{(2\pi f_c)^2L_1}
    =\frac{1}{(2\pi\cdot 3400)^2\cdot 10^{-3}}
    \approx \SI{2.19}{\micro F}.
\end{equation}

Para una respuesta lo más plana posible se tomó:

\begin{equation}
    Q=\frac{1}{\sqrt{2}}\approx 0.707.
\end{equation}

En un RLC serie:

\begin{equation}
    Q=\frac{\omega_0L}{R}
    \quad \Rightarrow \quad
    R=\frac{\omega_0L}{Q}.
\end{equation}

Entonces:

\begin{equation}
    R_1=\frac{2\pi\cdot 3400\cdot 10^{-3}}{0.707}\approx \SI{30.2}{\ohm}.
\end{equation}

Se adoptaban comercialmente $C_1=\SI{2.2}{\micro F}$ y $R_1=\SI{30}{\ohm}$.

La etapa pasaaltos RC-RC, con $R_3=R_4=R$ y $C_3=C_4=C$, tiene una transferencia aproximada:

\begin{equation}
    \Hpa(s)=\frac{(sRC)^2}{(sRC)^2+3sRC+1}.
\end{equation}

\begin{figure}[H]
    \centering
    \includegraphics[width=0.9\textwidth]{\figpath bode_opcion_a.pdf}
    \caption{Bode de la alternativa A: pasabajos RLC, buffer y pasaaltos RC-RC.}
    \label{fig:opcion_a_bode}
\end{figure}

Esta alternativa fue descartada porque no quedaba simétrica en términos de implementación: una etapa era resonante RLC y la otra era una red RC-RC cargada. En consecuencia, las pendientes, fase y sensibilidad a tolerancias no resultaban parejas (Figura \ref{fig:opcion_a_bode}).

\subsection{Alternativa B: dos etapas RLC con buffer}

La segunda alternativa consistía en utilizar dos etapas RLC, una pasabajos y una pasaaltos, separadas por un buffer. Esta opción era más simétrica desde el punto de vista conceptual.

\begin{figure}[H]
    \centering
\begin{tikzpicture}[transform shape]
	% Paths, nodes and wires:
	\draw (4.25, 1) to[sinusoidal current source, l={$v_i$}] (4.25, -1.25);
	\draw (4.25, 1) to[american resistor, l={$R_1$}] (6.25, 1);
	\draw (6.25, 1) to[american inductor, l={$L1$}] (8, 1);
	\draw (8, 1) to[capacitor, l={$C_1$}] (8, -1.25);
	\node[ground] at (8, -1.25){};
	\node[op amp, yscale=-1] at (10.333, 0.51){};
	\draw (8, 1) -| (9.143, 1.029);
	\draw (9.143, 0.02) |- (11.5, -1.75) -| (11.5, 0.5);
	\node[ground] at (4.25, -1.25){};
	\draw (12, 1) to[american resistor, l={$R_2$}] (13.5, 1);
	\draw (13.5, 1) to[capacitor, l={$C_2$}] (15, 1);
	\draw (15, 1) to[american inductor, l={$L_2$}] (15, -1.25);
	\node[ground] at (15, -1.25){};
	\draw (11.523, 0.51) |- (12, 1);
    \draw (15, 1) -- (15.75, 1);
    \node at (15.75, 1) [right] {$v_o$};
\end{tikzpicture}
\caption{Alternativa B: dos etapas RLC separadas por buffer.}
\label{fig:opcion_b}
\end{figure}

Para la etapa pasaaltos RLC con salida en el inductor:

\begin{equation}
    \Hpa(s)=\frac{s^2L_2C_2}{s^2L_2C_2+sR_2C_2+1}.
\end{equation}

El gráfico de Bode de la alternativa B se muestra en la Figura \ref{fig:opcion_b_bode}.

\begin{figure}[H]
    \centering
    \includegraphics[width=0.9\textwidth]{\figpath bode_opcion_b.pdf}
    \caption{Bode de la alternativa B: dos etapas RLC separadas por buffer.}
    \label{fig:opcion_b_bode}
\end{figure}

Si se fija $f_c=\SI{300}{Hz}$ y $L_2=\SI{1}{mH}$, el capacitor necesario es:

\begin{equation}
    C_2=\frac{1}{(2\pi\cdot 300)^2\cdot 10^{-3}}\approx \SI{281}{\micro F}.
\end{equation}

Con $Q=0.707$:

\begin{equation}
    R_2=\frac{2\pi\cdot 300\cdot 10^{-3}}{0.707}\approx \SI{2.66}{\ohm}.
\end{equation}

Este valor de resistencia es demasiado bajo para ser atacado directamente por la salida del TL082. En la simulación transitoria se obtuvo una corriente de salida del orden de:

\begin{equation}
    I_{\mathrm{opamp}}\approx\SI{27}{mA}.
\end{equation}

% \begin{figure}[H]
%     \centering
%     \includegraphics[width=0.82\textwidth]{\figpath verificacion_corriente_tl082.png}
%     \caption{Verificación de corriente exigida al TL082 en una alternativa RLC con $L=\SI{1}{mH}$.}
%     \label{fig:corriente_tl082}
% \end{figure}

Este resultado motivó descartar la alternativa RLC con inductores de \SI{1}{mH}. Matemáticamente cumplía la función, pero físicamente exigía demasiada corriente al operacional.

\subsection{Rediseño RLC con inductores comerciales mayores}

Para que las resistencias de las etapas RLC no fueran de pocos ohms, se analizó qué inductancia sería necesaria. Partiendo de:

\begin{equation}
    R=\frac{\omega_0L}{Q},
\end{equation}

se despeja:

\begin{equation}
    L=\frac{RQ}{\omega_0}.
\end{equation}

Para el corte inferior de \SI{300}{Hz}, si se busca $R\approx\SI{270}{\ohm}$:

\begin{equation}
    L_2=\frac{270\cdot0.707}{2\pi\cdot300}\approx\SI{101}{mH}.
\end{equation}

Se obtiene un valor comercial cercano de \SI{100}{mH}. El capacitor asociado sería:

\begin{equation}
    C_2=\frac{1}{(2\pi\cdot300)^2\cdot\SI{100}{mH}}\approx\SI{2.81}{\micro F},
\end{equation}

por lo que puede adoptarse $C_2=\SI{2.7}{\micro F}$.

Para el corte superior de \SI{3.4}{kHz}, si se busca $R\approx\SI{300}{\ohm}$:

\begin{equation}
    L_1=\frac{300\cdot0.707}{2\pi\cdot3400}\approx\SI{9.9}{mH}.
\end{equation}

Se obtiene un valor comercial de \SI{10}{mH}. El capacitor asociado es:

\begin{equation}
    C_1=\frac{1}{(2\pi\cdot3400)^2\cdot\SI{10}{mH}}\approx\SI{219}{nF},
\end{equation}

por lo que se adopta $C_1=\SI{220}{nF}$.

Esta alternativa solucionaba el problema de corriente del TL082, pero los inductores de \SI{100}{mH} comerciales tienen mucha resistencia serie, y junto a la tolerancia propia de los componentes la respuesta de Bode terminaba por degradarse. Por este motivo, se decidió descartar la alternativa RLC-RLC.

La respuesta de esta alternativa resulta igual a la anterior, pero con inductores más grandes y resistencias más altas. La Figura \ref{fig:bode_opcion_c} muestra la respuesta de Bode de esta alternativa.
\begin{figure}[H]
    \centering
    \includegraphics[width=0.86\textwidth]{\figpath bode_opcion_c.pdf}
    \caption{Bode filtro de dos etapas RLC con valores modificados.}
    \label{fig:bode_opcion_c}
\end{figure}

\subsection{Diseño final adoptado: RC-RC con buffers}

El diseño final utiliza un pasabajos RC-RC, un buffer, un pasaaltos CR-CR y un buffer de salida. La topología final es compatible con la restricción de la consigna, ya que el TL082 se utiliza como seguidor de tensión.

\begin{figure}[H]
    \centering
    \begin{tikzpicture}[transform shape]
	% Paths, nodes and wires:
	\draw (3.357, -1.029) to[sinusoidal current source, l={$v_i$}] (3.357, -3.279);
	\node[op amp, yscale=-1] at (10.44, -1.49){};
	\draw (9.273, -1.97) |- (11.63, -3.74) -| (11.63, -1.49);
	\node[ground] at (3.357, -3.279){};
	\draw (11.63, -1.49) |- (12.107, -1);
	\draw (3.357, -1.029) to[american resistor, l={$R_1$}] (5.5, -1);
	\draw (5.5, -1) to[american resistor, l={$R_2$}] (9.25, -1);
	\draw (5.5, -1) to[capacitor, l={$C_1$}] (5.5, -3.25);
	\draw (8.5, -1) to[capacitor, l_={$C_2$}] (8.5, -3.25);
	\draw (12.107, -1) to[capacitor, l={$C_3$}] (13.75, -1);
	\draw (13.75, -1) to[capacitor, l={$C_4$}] (15.75, -1);
	\draw (13.75, -1) to[american resistor, l={$R_3$}] (13.75, -3.25);
	\draw (15.75, -1) to[american resistor, l_={$R_4$}] (15.75, -3.25);
	\node[op amp, yscale=-1] at (17.69, -1.49){};
	\draw (16.5, -1.02) |- (15.75, -1);
	\draw (16.5, -1.99) -| (16.5, -2.99) -- (19, -2.99) |- (18.88, -1.5);
	\draw (19, -1.5) -- (19, -1) -- (20, -1);
	\node[ground] at (5.5, -3.25){};
	\node[ground] at (8.5, -3.25){};
	\node[ground] at (13.75, -3.25){};
	\node[ground] at (15.75, -3.25){};
    \node at (20, -1) [right] {$v_o$};
\end{tikzpicture}
    \caption{Filtro final adoptado: pasabajos RC-RC, buffer, pasaaltos RC-RC y buffer de salida.}
    \label{fig:filtro_final_circuitikz}
\end{figure}

\subsection{Cálculo de los valores RC}

Para dos etapas RC iguales en cascada sin buffer intermedio, el pasabajos queda:

\begin{equation}
    \Hpb(s)=\frac{1}{(sRC)^2+3sRC+1}.
\end{equation}

Esta red no tiene el corte de \SI{-3}{dB} exactamente en $1/(2\pi RC)$. Para esta topología, el punto de \SI{-3}{dB} aparece aproximadamente en:

\begin{equation}
    f_{c,PB}\approx 0.374\frac{1}{2\pi RC}.
\end{equation}

Se desea $f_{c,PB}\approx\SI{3.4}{kHz}$. Por lo tanto:

\begin{equation}
    \frac{1}{2\pi RC}\approx\frac{\SI{3.4}{kHz}}{0.374}\approx\SI{9.1}{kHz}.
\end{equation}

Adoptando $C_1=C_2=\SI{10}{nF}$:

\begin{equation}
    R_1=R_2\approx\frac{1}{2\pi\cdot \SI{9.1}{kHz}\cdot \SI{10}{nF}}\approx\SI{1.75}{k\ohm}.
\end{equation}

Se adopta el valor comercial:

\begin{equation}
    \boxed{R_1=R_2=\SI{1.8}{k\ohm}},
    \qquad
    \boxed{C_1=C_2=\SI{10}{nF}}.
\end{equation}

Para el pasaaltos RC-RC:

\begin{equation}
    \Hpa(s)=\frac{(sRC)^2}{(sRC)^2+3sRC+1}.
\end{equation}

El punto de \SI{-3}{dB} aparece aproximadamente en:

\begin{equation}
    f_{c,PA}\approx 2.67\frac{1}{2\pi RC}.
\end{equation}

Se desea $f_{c,PA}\approx\SI{300}{Hz}$. Por lo tanto:

\begin{equation}
    \frac{1}{2\pi RC}\approx\frac{\SI{300}{Hz}}{2.67}\approx\SI{112}{Hz}.
\end{equation}

Adoptando $C_3=C_4=\SI{100}{nF}$:

\begin{equation}
    R_3=R_4\approx\frac{1}{2\pi\cdot\SI{112}{Hz}\cdot\SI{100}{nF}}\approx\SI{14.2}{k\ohm}.
\end{equation}

Se adopta el valor comercial:

\begin{equation}
    \boxed{R_3=R_4=\SI{15}{k\ohm}},
    \qquad
    \boxed{C_3=C_4=\SI{100}{nF}}.
\end{equation}

\subsection{Transferencia total del filtro final}

La transferencia total idealizada se obtiene como:

\begin{equation}
    H(s)=\Hpb(s)\Hpa(s).
\end{equation}

Reemplazando las expresiones:

\begin{equation}
    \boxed{
    H(s)=
    \frac{1}{(sR_{PB}C_{PB})^2+3sR_{PB}C_{PB}+1}
    \cdot
    \frac{(sR_{PA}C_{PA})^2}{(sR_{PA}C_{PA})^2+3sR_{PA}C_{PA}+1}
    }
\end{equation}

con:

\begin{equation}
    R_{PB}=\SI{1.8}{k\ohm}, \quad C_{PB}=\SI{10}{nF}, \quad
    R_{PA}=\SI{15}{k\ohm}, \quad C_{PA}=\SI{100}{nF}.
\end{equation}

El sistema tiene dos ceros en el origen aportados por la etapa pasaaltos. Los polos provienen de los dos denominadores cuadráticos de las redes RC-RC. La ubicación de estos polos determina la pendiente y el ancho de la banda útil.
Los polos de las etapas combinadas resultan:
\begin{enumerate}
    \item Polos de la etapa pasabajos RC-RC:
    \begin{equation}
        p_{PB,1}\approx -145446.33\text{ rad/s},
        \qquad
        p_{PB,2}\approx -21220.33\text{ rad/s}.
    \end{equation}
    \item Polos de la etapa pasaaltos RC-RC:
    \begin{equation}
        p_{PA,1}\approx -1745.36\text{ rad/s},
        \qquad
        p_{PA,2}\approx -254.64\text{ rad/s}.
    \end{equation}
\end{enumerate}

% Ceros:
% z1, 2 = 0 (cero doble en el origen)
% Polos:
% p1 = -145446.33 rad/s
% p2 = -21220.33 rad/s
% p3 = -1745.36 rad/s
% p4 = -254.64 rad/s


Los polos y ceros del filtro final se muestran en la Figura \ref{fig:polos_ceros_filtro_final}. La respuesta en frecuencia se muestra en la Figura \ref{fig:bode_filtro_final}.


\begin{figure}[H]
    \centering
    \includegraphics[width=0.86\textwidth]{\figpath diagrama_polos_ceros_filtro_final.pdf}
    \caption{Diagrama de polos y ceros del filtro final.}
    \label{fig:polos_ceros_filtro_final}
\end{figure}

\subsection{Valores finales}

\begin{table}[H]
    \centering
    \begin{tabular}{ccc}
        \toprule
        Etapa & Componente & Valor adoptado \\
        \midrule
        Pasabajos & $R_1=R_2$ & \SI{1.8}{k\ohm} \\
        Pasabajos & $C_1=C_2$ & \SI{10}{nF} \\
        Pasaaltos & $R_3=R_4$ & \SI{15}{k\ohm} \\
        Pasaaltos & $C_3=C_4$ & \SI{100}{nF} \\
        Buffers & U1A, U1B & TL082 \\
        Alimentación & $+V$, $-V$ & $\pm\SI{12}{V}$ \\
        \bottomrule
    \end{tabular}
    \caption{Valores adoptados para el filtro final.}
    \label{tab:valores_filtro_final}
\end{table}

\subsection{Esquemático en Altium}

El diseño se implementó en Altium con entrada por conector, entrada por jack, salida por conector y bornera de alimentación. Se incluyeron dos seguidores de tensión del TL082, uno para desacoplar las etapas y otro para entregar la señal de salida.

\begin{figure}[H]
    \centering
    \includegraphics[width=0.98\textwidth]{\figpath esquematico_altium.png}
    \caption{Esquemático final implementado en Altium.}
    \label{fig:esquematico_altium}
\end{figure}
